% ****** Start of file aipsamp.tex ******
%
% This file follows the REVTeX/AIP template used in mainv11.tex.
\documentclass[%
 aip,
 jmp,%
 amsmath,amssymb,
 reprint,%
]{revtex4-2}

\usepackage{graphicx}
\usepackage{dcolumn}
\usepackage{bm}
\usepackage{mathtools}
\usepackage{physics}
\usepackage{amsthm}
\usepackage{xcolor}
\usepackage{hyperref}

\begin{document}

\preprint{AIP/123-QED}

\title[Bogoliubov transformation]{Bogoliubov transformation:\\
Remaining analytic solutions for $\zeta_1(x)$}

\author{Javier Huenupi}
\homepage{https://www.cec.uchile.cl/~javier.huenupi/}
\email{javier.huenupi@ug.uchile.cl}
\author{Gonzalo Palma}
\affiliation{
Departamento de F\'isica, FCFM, Universidad de Chile, Blanco Encalada 2008,
Santiago, Chile
}%
\author{Claudio Mu\~noz}%
\affiliation{%
Departamento de Ingenier\'ia Matem\'atica and Centro de Modelamiento Matem\'atico
(UMI 2807 CNRS), Universidad de Chile, Casilla 170 Correo 3, Santiago, Chile.
}%

\date{\today}

\begin{abstract}
We show that the fourth-order equation for the Bogoliubov coefficient
$\zeta_1(x)$ hides a second-order Kummer equation.  The two solutions already
obtained by the Laplace-transform method correspond to the Kummer
$\,_1F_1$ branch.  The two remaining solutions are obtained by applying the
same differential map to the Tricomi branch $U(a,b,z)$.
\end{abstract}

\keywords{inflation, Bogoliubov, multi-field inflation}
\maketitle

\newcommand{\z}{\zeta}
\newcommand{\p}{\psi}
\newcommand{\K}{\mathcal{K}}
\newcommand{\M}{M}
\newcommand{\U}{U}

\section{The equation to solve}

The main equation in the previous document is
\begin{equation}
\label{eq:xode}
\begin{aligned}
0 ={}& x^2 \zeta_1^{(4)}(x)
 + 4x(1+ix)\zeta_1^{(3)}(x)
 + 2x(5i-2x)\zeta_1''(x)  \\
& - 2(i+2x)\zeta_1'(x)
 - \lambda^2 \zeta_1(x),
\qquad x>0,\quad \lambda>0.
\end{aligned}
\end{equation}
The Laplace-transform method previously used gives two linearly independent
solutions that are bounded at $x=0$.  The goal is to find the other two
solutions of the fourth-order equation.

The useful observation is that Eq.~\eqref{eq:xode} is not an arbitrary
fourth-order equation.  After a simple complex change of variable it can be
generated from Kummer's second-order equation.

\section{Change of variable}

Let
\begin{equation}
\label{eq:zdef}
z=-2ix,\qquad \zeta_1(x)=Y(z).
\end{equation}
Since $d/dx=-2i\,d/dz$, Eq.~\eqref{eq:xode} is equivalent, after division by
$-4$, to
\begin{equation}
\label{eq:zode}
z^2Y^{(4)}
+ (4z-2z^2)Y^{(3)}
+ (z^2-5z)Y''
+ (1+z)Y'
+ \frac{\lambda^2}{4}Y=0.
\end{equation}
This form is the key simplification.

\section{Hidden Kummer equation}

Consider Kummer's equation
\begin{equation}
\label{eq:kummer}
\K_q[w]\equiv z w''+(1-z)w'-q w=0.
\end{equation}
Its two standard independent solutions are
\begin{equation}
w(z)=\,_1F_1(q;1;z),
\qquad
w(z)=U(q,1,z),
\end{equation}
where $U$ is Tricomi's confluent hypergeometric function.

Now define the first-order map
\begin{equation}
\label{eq:map}
Y(z)=(1+2q)w(z)-2w'(z).
\end{equation}
If $w$ solves Eq.~\eqref{eq:kummer}, then all higher derivatives of $w$ can be
reduced to a linear combination of $w$ and $w'$.  Substitution into the
left-hand side of Eq.~\eqref{eq:zode} gives
\begin{equation}
\begin{multlined}
z^2Y^{(4)}
+(4z-2z^2)Y^{(3)}
+(z^2-5z)Y''
+(1+z)Y'
+\frac{\lambda^2}{4}Y \\
=
\left(\frac{\lambda^2}{4}+q^2\right)Y.
\end{multlined}
\end{equation}
Therefore, if
\begin{equation}
\label{eq:qvalues}
q=\pm \frac{i\lambda}{2},
\end{equation}
then $q^2=-\lambda^2/4$, and the right-hand side vanishes.  This proves that
Eq.~\eqref{eq:map} sends every solution of Kummer's equation into a solution
of the original fourth-order equation.

\section{Recovering the two known solutions}

First choose the Kummer $\,_1F_1$ branch:
\begin{equation}
w(z)=\,_1F_1(q;1;z).
\end{equation}
Using
\begin{equation}
\dv{}{z}\,_1F_1(q;1;z)=q\,_1F_1(q+1;2;z),
\end{equation}
the map \eqref{eq:map} gives
\begin{equation}
\label{eq:Mbranch}
Y_M(z;q)
=(1+2q)\,_1F_1(q;1;z)
-2q\,_1F_1(q+1;2;z).
\end{equation}

For $q=-i\lambda/2$ and $z=-2ix$,
\begin{equation}
\label{eq:known1bare}
\begin{multlined}
Y_M\left(-2ix;-\frac{i\lambda}{2}\right)
=(1-i\lambda)\,_1F_1\left(-\frac{i\lambda}{2};1;-2ix\right)\\
+i\lambda\,_1F_1\left(1-\frac{i\lambda}{2};2;-2ix\right).
\end{multlined}
\end{equation}
Since
\begin{equation}
L_{\nu}(z)=\,_1F_1(-\nu;1;z),
\end{equation}
this is precisely the structure of the previous solution
\begin{equation}
\label{eq:zeta11}
\begin{multlined}
\zeta_{1,1}(x)
=f_0(\lambda)
\Bigg[
(1-i\lambda)L_{i\lambda/2}(-2ix)\\
+i\lambda\,_1F_1\left(1-\frac{i\lambda}{2};2;-2ix\right)
\Bigg].
\end{multlined}
\end{equation}
Similarly, for $q=+i\lambda/2$,
\begin{equation}
\label{eq:zeta12}
\begin{multlined}
\zeta_{1,2}(x)
=f_0(-\lambda)
\Bigg[
(1+i\lambda)L_{-i\lambda/2}(-2ix)\\
-i\lambda\,_1F_1\left(1+\frac{i\lambda}{2};2;-2ix\right)
\Bigg].
\end{multlined}
\end{equation}
Thus the known bounded pair is simply the $\,_1F_1$ branch of the hidden
Kummer problem.

\section{The missing pair}

The remaining branch of Kummer's equation is Tricomi's function
\begin{equation}
w(z)=U(q,1,z).
\end{equation}
The derivative identity is
\begin{equation}
\dv{}{z}U(q,1,z)=-q\,U(q+1,2,z).
\end{equation}
Therefore, the same map \eqref{eq:map} gives
\begin{equation}
\label{eq:Ubranch}
Y_U(z;q)
=(1+2q)U(q,1,z)
+2q\,U(q+1,2,z).
\end{equation}
Taking again $q=\pm i\lambda/2$ and $z=-2ix$, a convenient unnormalized
choice of the two remaining solutions is
\begin{equation}
\label{eq:zeta13}
\begin{multlined}
\zeta_{1,3}(x)
=(1-i\lambda)U\left(-\frac{i\lambda}{2},1,-2ix\right)\\
-i\lambda U\left(1-\frac{i\lambda}{2},2,-2ix\right),
\end{multlined}
\end{equation}
and
\begin{equation}
\label{eq:zeta14}
\begin{multlined}
\zeta_{1,4}(x)
=(1+i\lambda)U\left(\frac{i\lambda}{2},1,-2ix\right)\\
+i\lambda U\left(1+\frac{i\lambda}{2},2,-2ix\right).
\end{multlined}
\end{equation}
Multiplying either function by a non-zero constant gives the same solution
space.  Also, adding any linear combination of $\zeta_{1,1}$ and
$\zeta_{1,2}$ to either of these functions gives an equivalent basis.

\section{Why these solutions were missed}

The Laplace-transform construction imposed boundedness at $x=0$ and used the
relation
\begin{equation}
-2i\zeta_1'(0)=\lambda^2\zeta_1(0).
\end{equation}
The $U$ branch does not satisfy this boundedness hypothesis.  Indeed, for
$b=2$, $U(q+1,2,z)$ has a leading $1/z$ singularity at $z=0$, and since
$z=-2ix$, the terms in Eqs.~\eqref{eq:zeta13}--\eqref{eq:zeta14} contain a
singular contribution of order $1/x$ plus logarithmic terms.  This matches the
local analysis around $x=0$, where one expects singular solutions such as
\begin{equation}
\zeta_1(x)\sim \frac{6}{x}+6i\log x+\cdots.
\end{equation}
Thus the $U$ branch supplies precisely the missing singular sector.

\section{Relation with the large-$x$ basis}

The four functions
\begin{equation}
Y_M(z;-i\lambda/2),\quad
Y_M(z;i\lambda/2),\quad
Y_U(z;-i\lambda/2),\quad
Y_U(z;i\lambda/2)
\end{equation}
form a basis of the fourth-order equation.  The $U$ branch is not necessarily
the most convenient basis if one wants the asymptotic modes normalized exactly
as
\begin{equation}
x^{-1+i\lambda/2},\qquad x^{-1-i\lambda/2}.
\end{equation}
This is because $U(q,1,z)\sim z^{-q}$ as $|z|\to\infty$, while the
$\,_1F_1$ branch contains both an algebraic part and an exponential part.

If a subleading Jost-type basis at infinity is desired, one may cancel the
leading algebraic piece by forming
\begin{equation}
\label{eq:jostcomb}
Y_{\rm sub}(z;q)
=Y_M(z;q)
-\frac{e^{i\pi q}}{\Gamma(1-q)}Y_U(z;q),
\end{equation}
up to branch conventions for $\arg z$.  This isolates the
$e^z z^{q-1}$ part of the Kummer asymptotics.  Since $e^z=e^{-2ix}$ and
$q=\pm i\lambda/2$, these combinations are the natural candidates for the
oscillatory large-$x$ branches.

\section{Recovering $\zeta_2,\psi_1,\psi_2$ from $\zeta_1$}

Once $\zeta_1$ is known, the other three Bogoliubov coefficients do not need to
be obtained by solving new fourth-order equations.  They can be reconstructed
algebraically from the first-order system.  In the $x=e^{-t}$ variable this
system is
\begin{equation}
\label{eq:firstordersystem}
\begin{aligned}
\zeta_1'
&=
\frac{i\lambda}{2x^2}
\left[(ix-1)\psi_1+(ix+1)e^{-2ix}\psi_2\right],\\
\zeta_2'
&=
\frac{i\lambda}{2x^2}
\left[(ix-1)e^{2ix}\psi_1+(ix+1)\psi_2\right],\\
\psi_1'
&=
\frac{i\lambda}{2x^2}
\left[-(ix+1)\zeta_1+(ix+1)e^{-2ix}\zeta_2\right],\\
\psi_2'
&=
\frac{i\lambda}{2x^2}
\left[(ix-1)e^{2ix}\zeta_1-(ix-1)\zeta_2\right].
\end{aligned}
\end{equation}

Let
\begin{equation}
y(x):=\zeta_1(x).
\end{equation}
We first use the first equation of \eqref{eq:firstordersystem}.  Write it as
\begin{equation}
\label{eq:yprime}
y'=A\psi_1+B\psi_2,
\end{equation}
where
\begin{equation}
A=\frac{i\lambda}{2x^2}(ix-1),
\qquad
B=\frac{i\lambda}{2x^2}(ix+1)e^{-2ix}.
\end{equation}
Differentiating Eq.~\eqref{eq:yprime} and using the equations for
$\psi_1'$ and $\psi_2'$ in \eqref{eq:firstordersystem}, the terms proportional
to $\zeta_2$ cancel.  One obtains
\begin{equation}
\label{eq:ysecond}
y''
=
\frac{i\lambda e^{-2ix}}{2x^3}
\left[
(-ix+2)e^{2ix}\psi_1
+(2x^2-3ix-2)\psi_2
\right].
\end{equation}
Therefore Eqs.~\eqref{eq:yprime} and \eqref{eq:ysecond} are two algebraic
equations for $\psi_1$ and $\psi_2$.  Solving them gives
\begin{equation}
\label{eq:psi1fromy}
\psi_1
=
\frac{
(ix^2+x)y''
+(-2x^2+3ix+2)y'
}
{\lambda x},
\end{equation}
and
\begin{equation}
\label{eq:psi2fromy}
\psi_2
=
-\frac{i e^{2ix}}{\lambda x}
\left[
(x^2+ix)y''
+(x+2i)y'
\right].
\end{equation}

To recover $\zeta_2$, differentiate Eq.~\eqref{eq:ysecond} once more and again
use the first-order system to remove $\psi_1'$, $\psi_2'$, and $\zeta_2'$.
Then substitute Eqs.~\eqref{eq:psi1fromy}--\eqref{eq:psi2fromy}.  The resulting
equation is algebraic in $\zeta_2$ and gives
\begin{equation}
\label{eq:zeta2fromy}
\begin{multlined}
\zeta_2
=
\frac{e^{2ix}}{\lambda^2}
\Big[
\lambda^2 y
-2ix^2y^{(3)}
+4x(x-i)y''\\
+4(x+i)y'
\Big].
\end{multlined}
\end{equation}

Thus any solution $y=\zeta_1$ of Eq.~\eqref{eq:xode} generates a solution of
the full first-order system by
\begin{equation}
\label{eq:fullmap}
y
\longmapsto
\left(
\zeta_1,\zeta_2,\psi_1,\psi_2
\right)
=
\left(
y,\zeta_2[y],\psi_1[y],\psi_2[y]
\right),
\end{equation}
where $\psi_1[y]$, $\psi_2[y]$, and $\zeta_2[y]$ are given by
Eqs.~\eqref{eq:psi1fromy}, \eqref{eq:psi2fromy}, and
\eqref{eq:zeta2fromy}.

The consistency check is direct.  Substituting the reconstruction formulas into
the equations for $\zeta_1'$, $\psi_1'$, and $\psi_2'$ gives identities.  The
remaining equation for $\zeta_2'$ gives
\begin{equation}
\label{eq:zeta2residual}
\begin{multlined}
\zeta_2'
-
\frac{i\lambda}{2x^2}
\left[(ix-1)e^{2ix}\psi_1+(ix+1)\psi_2\right]\\
=
-\frac{2ie^{2ix}}{\lambda^2}
\left[
x^2y^{(4)}
+4x(1+ix)y^{(3)}
+2x(5i-2x)y''
-2(i+2x)y'
-\lambda^2 y
\right].
\end{multlined}
\end{equation}
The bracket on the right-hand side is exactly the fourth-order equation
\eqref{eq:xode} for $y=\zeta_1$.  Hence the reconstructed functions solve the
full coupled system whenever $y$ solves the $\zeta_1$ equation.

\section{Checks}

There are three independent checks.

First, the derivation itself proves that Eqs.~\eqref{eq:zeta13} and
\eqref{eq:zeta14} solve Eq.~\eqref{eq:xode}: they are obtained from solutions
of Kummer's equation through a map that sends Kummer solutions into solutions
of Eq.~\eqref{eq:zode}.

Second, direct substitution of Eqs.~\eqref{eq:zeta13} and \eqref{eq:zeta14}
into Eq.~\eqref{eq:xode} gives zero.  In Mathematica, this check can be
performed with
\begin{verbatim}
odeX[zeta13Candidate] // FullSimplify
odeX[zeta14Candidate] // FullSimplify
\end{verbatim}
or numerically at high precision if the symbolic hypergeometric identities do
not simplify automatically.

Third, the Wronskian of the four functions
\begin{equation}
\{\zeta_{1,1},\zeta_{1,2},\zeta_{1,3},\zeta_{1,4}\}
\end{equation}
is non-zero at a regular point, confirming linear independence.

\section{Conclusion}

The main insight is that the fourth-order equation for $\zeta_1$ is generated
from a second-order Kummer equation.  The two known solutions are the
$\,_1F_1$ branch.  The two missing solutions are the $U$ branch:
\begin{multline*}
\zeta_{1,3}(x)
=(1-i\lambda)U\left(-\frac{i\lambda}{2},1,-2ix\right)\\
-i\lambda U\left(1-\frac{i\lambda}{2},2,-2ix\right),
\end{multline*}
\begin{multline*}
\zeta_{1,4}(x)
=(1+i\lambda)U\left(\frac{i\lambda}{2},1,-2ix\right)\\
+i\lambda U\left(1+\frac{i\lambda}{2},2,-2ix\right).
\end{multline*}
These functions complete the four-dimensional solution space.  Their
singularity at $x=0$ explains why they were not obtained by the bounded
Laplace-transform method.

\nocite{*}
\bibliography{bibmain}

\end{document}
